% section_3.tex — Evaluation and Results
% Current participant counts: cross-method n=18, convergence n=12.

\section{Evaluation and Results}
\label{sec:results}

\subsection{Definition of success}
\label{sec:success}

A useful evaluation harness has to satisfy two parties, as flagged at
the end of \S\ref{sec:background}. We restate the two criteria here in
concrete, checkable form.

\paragraph{Researcher-side success.}
The harness should give a graphics researcher data they can act on
without engineering the data pipeline themselves. Concretely:
(a)~every participant who starts a session leaves a snapshot JSON on
disk, including sessions that ended early; (b)~the resulting data has
informative structure --- within-participant slider spread is
non-trivial, A/B preferences are not all ties, and aspect citations
are distributed across the three options rather than degenerate on one;
and (c)~the data shape is general enough that a researcher can apply
analyses of their own choosing (correlations against offline metrics,
within-trial regressions, inter-rater agreement) without bespoke
parsing.

\paragraph{Participant-side success.}
The session should feel smooth and easy to complete. Concretely:
(a)~interactive latency is below the threshold at which participants
notice it as ``laggy''; (b)~both renderers sustain their target frame
rate over the full session without sustained dips; (c)~no participant
abandons mid-session for technical reasons; and (d)~the rating UI is
clear enough that no experimenter intervention is required.

Participant-side success is harder to measure than researcher-side
success because it is partly experiential: systems telemetry can
confirm that no frames were dropped, but it cannot confirm that the
participant felt the session was pleasant. We rely on telemetry as an
indirect proxy for now, and plan to administer a brief post-session
self-report survey (in preparation) covering perceived smoothness,
clarity of the rating UI, and fatigue (\S\ref{sec:participant_feedback}).
The dual-party framing also clarifies what kind of failure invalidates
what kind of conclusion: a participant-side failure (frame stalls, UI
confusion) propagates into researcher-side measurements and invalidates
them, while a researcher-side data-shape ``failure'' --- sliders all
pegged at one end, for instance --- could itself be a valid measurement
of a real ceiling effect rather than a harness bug.

\subsection{Case studies}
\label{sec:setup}

We exercised the harness with two case studies. Both ran on a single
AWS GPU instance with an NVIDIA L4 (24~GB VRAM, \texttt{g6.xlarge}),
hosting all server-side processes (\texttt{mediamtx}, \texttt{ws\_mux},
\texttt{serve.py}, and the active renderers) on the same host.
Participants connected from their own machines over the public
internet; we did not constrain network conditions or display
characteristics, and we discuss the resulting variance in
\S\ref{sec:limitations}. Across both studies we recruited 19 unique
participants from the authors' friends and labmates; eleven of the
twelve case-B participants had previously completed case A. A case-A session
lasted a median of 9~minutes (mean 10); a case-B session a median of
5~minutes. We collected $n=18$ complete sessions in case study A and
$n=12$ in case study B; recruitment is closed for this submission.

\paragraph{Case study A: comparison (i-NGP vs Nexels).}
Two NVS methods --- i-NGP and Nexels --- rendered side-by-side at a
matched storage budget ($\approx$14~MB per scene; see
Table~\ref{tab:cross_offline}) over six mip-NeRF-360 scenes
(\texttt{bicycle}, \texttt{counter}, \texttt{garden}, \texttt{kitchen},
\texttt{room}, \texttt{stump}) plus one practice scene
(\texttt{bonsai}). The two methods stream concurrently to two
side-by-side viewports at 1440$\times$810@30. Method identity is
hidden on the viewport badges (only the \texttt{test} initials hook
reveals it, and is never used in real sessions). The seven scene pairs
are presented in a per-participant random order; side assignment of
i-NGP and Nexels is coin-flipped per pair; the first pair
(\texttt{bonsai}) is the practice trial and is excluded from analysis.
The questionnaire shape --- per-side sliders, forced A/B, aspect
multi-select --- is described in \S\ref{sec:questionnaire}.

\paragraph{Case study B: convergence (3DGS at four checkpoint counts).}
A single method (3D Gaussian Splatting) trained for varying numbers of
iterations (1k, 3k, 7k, 30k) on two Tanks-and-Temples scenes
(\texttt{train}, \texttt{truck}) --- eight stimuli total. All eight
checkpoints are GPU-resident ($\approx$1.4~GB combined, 210k--3.8M
splats per checkpoint), so swapping the active checkpoint is a
constant-time pointer update. The renderer streams 1920$\times$1080@60.
Before each scene's rating block, the participant watches a
low-frame-rate slideshow of the training photographs as a
ground-truth reference. Scene order is randomized; iteration order
within a scene is fixed ascending so the reference frame for
\emph{better} stays stable across the four levels.

\paragraph{Offline-metric baselines.}
For each scene--method pair in case study A and each
scene--checkpoint pair in case study B we computed PSNR, SSIM, and
LPIPS on the held-out test split for that scene (mip-NeRF-360
hold-every-eighth-view convention for case study A; gsplat
validation split for case study B). LPIPS uses the AlexNet backbone in
both cases for consistency. Tables~\ref{tab:cross_offline}
and~\ref{tab:convergence} list these per-stimulus offline numbers
alongside the corresponding mean slider scores.

\begin{table}[t]
\begin{minipage}[t]{0.48\textwidth}
\centering\footnotesize
\captionof{table}{Case study A: \emph{offline metrics} at the
\texttt{small} size setting ($\approx$14~MB per scene). PSNR / SSIM /
LPIPS are computed on the held-out mip-NeRF-360 test split with a
consistent LPIPS-AlexNet backbone; FPS is from each method's bench
script over 200 cycled test-camera frames at 1297$\times$840. Bold
marks the per-scene winner on each column. \texttt{bonsai} is the
practice scene and is excluded from the mean.}
\label{tab:cross_offline}
\setlength{\tabcolsep}{3pt}
\begin{tabular}{llrrrrr}
\toprule
Scene & Method & PSNR$\uparrow$ & SSIM$\uparrow$ & LPIPS$\downarrow$ & FPS & MB \\
\midrule
bicycle & i-NGP  & 21.93          & 0.462          & 0.592          & \textbf{62.0} & 14.5          \\
        & Nexels & \textbf{22.21} & \textbf{0.539} & \textbf{0.424} & 29.5          & \textbf{14.4} \\
counter & i-NGP  & 25.31          & 0.757          & 0.272          & \textbf{61.5} & \textbf{13.5} \\
        & Nexels & \textbf{26.42} & \textbf{0.833} & \textbf{0.216} & 30.0          & 14.4          \\
garden  & i-NGP  & 23.13          & 0.514          & 0.456          & \textbf{74.9} & \textbf{12.0} \\
        & Nexels & \textbf{23.46} & \textbf{0.636} & \textbf{0.322} & 34.0          & 14.4          \\
kitchen & i-NGP  & 27.12          & 0.767          & 0.214          & \textbf{96.2} & \textbf{12.6} \\
        & Nexels & \textbf{27.68} & \textbf{0.861} & \textbf{0.167} & 31.9          & 14.4          \\
room    & i-NGP  & 28.47          & 0.844          & 0.225          & \textbf{69.8} & 15.0          \\
        & Nexels & \textbf{29.15} & \textbf{0.877} & \textbf{0.207} & 29.8          & \textbf{14.4} \\
stump   & i-NGP  & 21.68          & 0.473          & 0.574          & 27.0          & 16.0          \\
        & Nexels & \textbf{23.88} & \textbf{0.604} & \textbf{0.398} & \textbf{32.9} & \textbf{14.4} \\
\midrule
\multirow{2}{*}{\textbf{mean}}
        & i-NGP  & 24.61          & 0.636          & 0.389          & \textbf{65.2} & 13.9          \\
        & Nexels & \textbf{25.47} & \textbf{0.725} & \textbf{0.289} & 31.4          & 14.4          \\
\bottomrule
\end{tabular}
\end{minipage}
\hfill
\begin{minipage}[t]{0.48\textwidth}
\centering\footnotesize
\captionof{table}{Case study A: \emph{per-scene human ratings}
($n=18$ participants). \textbf{fg}, \textbf{bg}, \textbf{sm} are mean
0--100 slider scores for foreground, background, and rendering
smoothness. \textbf{Votes} is the number of participants who preferred
that method on that scene; tie counts are bicycle 3, counter 1,
garden 4, kitchen 2, room 1, stump 1. Bold marks the per-scene winner
on each column.}
\label{tab:cross_human}
\setlength{\tabcolsep}{4pt}
\begin{tabular}{llrrrr}
\toprule
Scene & Method & fg$\uparrow$ & bg$\uparrow$ & sm$\uparrow$ & Votes \\
\midrule
bicycle & i-NGP  & 73.3          & 46.4          & 66.6          & 2              \\
        & Nexels & \textbf{80.6} & \textbf{68.3} & \textbf{72.5} & \textbf{13}    \\
counter & i-NGP  & \textbf{68.6} & 54.9          & \textbf{55.4} & 6              \\
        & Nexels & 67.9          & \textbf{59.2} & 48.1          & \textbf{10}    \\
garden  & i-NGP  & 71.3          & 62.3          & \textbf{56.9} & 7              \\
        & Nexels & \textbf{76.8} & \textbf{74.3} & 45.9          & 7              \\
kitchen & i-NGP  & \textbf{77.6} & \textbf{62.3} & \textbf{66.4} & \textbf{14}    \\
        & Nexels & 58.4          & 43.1          & 38.6          & 2              \\
room    & i-NGP  & \textbf{58.6} & \textbf{59.2} & \textbf{57.5} & \textbf{10}    \\
        & Nexels & 52.8          & 54.6          & 55.8          & 6              \\
stump   & i-NGP  & 60.7          & 52.9          & \textbf{63.9} & 5              \\
        & Nexels & \textbf{76.5} & \textbf{53.7} & 59.4          & \textbf{11}    \\
\midrule
\multirow{2}{*}{\textbf{mean}}
        & i-NGP  & 68.4          & 56.3          & \textbf{61.1} & 44             \\
        & Nexels & \textbf{68.9} & \textbf{58.9} & 53.4          & \textbf{49}    \\
\bottomrule
\end{tabular}
\end{minipage}
\end{table}

\begin{table}[t]
\centering
\caption{Case study B. Offline metrics and per-stimulus mean human
ratings ($n=12$) for the convergence-study stimuli (3DGS at four
checkpoint counts on Tanks-and-Temples \texttt{train} and
\texttt{truck}). LPIPS is AlexNet (gsplat default), internally
consistent within this table. FPS is \texttt{ellipse\_time} at the
gsplat eval resolution and is therefore an upper bound on the
interactive frame rate at 1920$\times$1080.}
\label{tab:convergence}
\small
\begin{tabular}{llrrrrrrrrr}
\toprule
Scene & Iters & PSNR$\uparrow$ & SSIM$\uparrow$ & LPIPS$\downarrow$ & FPS & \# Splats & fg$\uparrow$ & bg$\uparrow$ & sm$\uparrow$ \\
\midrule
train & 1k  & 17.06 & 0.584 & 0.550 & 433.7 & 0.21M & 48.8 & 30.9 & 72.5 \\
      & 3k  & 18.75 & 0.650 & 0.413 & 254.1 & 0.49M & 63.7 & 41.3 & 71.2 \\
      & 7k  & 19.73 & 0.721 & 0.298 & 130.8 & 1.05M & 78.3 & 51.1 & 80.2 \\
      & 30k & 21.44 & 0.812 & 0.146 &  77.4 & 1.81M & 84.2 & 64.2 & 79.2 \\
truck & 1k  & 20.94 & 0.725 & 0.326 & 437.5 & 0.21M & 43.2 & 48.0 & 66.8 \\
      & 3k  & 22.65 & 0.800 & 0.207 & 144.0 & 1.14M & 58.4 & 66.2 & 82.8 \\
      & 7k  & 23.87 & 0.845 & 0.143 &  67.6 & 2.51M & 77.5 & 72.0 & 82.9 \\
      & 30k & 25.04 & 0.871 & 0.095 &  45.5 & 3.79M & 81.3 & 71.8 & 60.1 \\
\bottomrule
\end{tabular}
\end{table}

\subsection{Systems results}
\label{sec:systems_results}

The harness met its interactivity targets in every session collected
to date. Both renderers in case study A sustained the target 30~fps
at 1440$\times$810 throughout; the 3DGS renderer in case study B
sustained 60~fps at 1920$\times$1080 across all four checkpoint
levels; no session was abandoned due to a renderer stall, dropped
WebRTC stream, or submission failure; and checkpoint swaps in case
study B introduced no measurable visual hitch. The per-method FPS
columns of Tables~\ref{tab:cross_offline} and~\ref{tab:convergence}
give the single-renderer ceilings at the offline benchmark resolution
(1297$\times$840 / gsplat eval). Per-scene FPS detail and the
end-to-end latency budget are deferred to Appendix~\ref{app:systems}.

\subsection{Data informativeness}
\label{sec:informativeness}

This subsection addresses the researcher-side success criterion: does
the harness produce data with the shape and signal required for
downstream analysis? We deliberately keep the discussion descriptive
rather than inferential. The substantive claims that one could make
\emph{from} this data --- about i-NGP vs Nexels, about the perceptual
return on additional 3DGS training, about how PSNR tracks human
preference --- are properly the researcher's to make. Our role here is
to verify that the harness has not produced degenerate data.

Three observations support the informativeness claim:

First, the slider distributions are spread.
Table~\ref{tab:cross_human} shows mean foreground/background/
smoothness scores in case study A ranging from 38 to 81 across
scene--method pairs, with no floor or ceiling effect.
Convergence-case slider means (Table~\ref{tab:convergence}) similarly
span 31 to 84 across the eight stimuli. Within-participant variance,
while not shown in the tables, is non-trivial: participants regularly
assigned different scores to different aspects of the same condition.

Second, A/B preferences in case study A are heterogeneous across
scenes (Figure~\ref{fig:preference_per_scene}). The aggregate
44/49/12 (i-NGP / Nexels / tie) is close to even, but the per-scene
distribution is anything but: \texttt{bicycle} swings 13--3--2 toward
Nexels (sign test $p=0.007$); \texttt{kitchen} swings 14--2--2 toward
i-NGP ($p=0.004$); the remaining four scenes sit closer to the
aggregate average. The data discriminates strongly between scenes ---
exactly the kind of structure that lets a downstream analysis say
something interesting.

\begin{figure}[h]
  \centering
  \includegraphics[width=0.85\linewidth]{fig_preference_per_scene}
  \caption{Per-scene A/B preference distribution from case study A
  ($n=18$). Each bar segments the 7-point preference scale into
  strong / moderate / slight i-NGP (red shades), tie (gray), and
  slight / moderate / strong Nexels (blue shades), so the bar's
  colour spread encodes preference \emph{strength} in addition to
  direction. Scenes ordered by net Nexels share, ascending. The
  cross-scene heterogeneity is the indication that the harness
  produces analyzable data; we do not interpret the per-scene
  rankings here.}
  \label{fig:preference_per_scene}
\end{figure}

Third, convergence-case slider trajectories across iteration counts
show non-trivial patterns (Figure~\ref{fig:convergence}). On the
smaller scene (\texttt{train}), foreground and background ratings rise
monotonically across all four checkpoints (foreground 49 $\to$ 84;
background 31 $\to$ 64), with smoothness flat after 7k. On
\texttt{truck}, the picture is more interesting: foreground saturates
between 7k and 30k (77 $\to$ 81), background flattens (72 $\to$ 72),
and \emph{smoothness drops 23 points} from 82.9 at 7k to 60.1 at 30k
--- a non-monotone trajectory that the data captures cleanly even at
$n=12$. The four checkpoint levels are clearly distinguishable, and
the data admits straightforward visual comparison against the offline
LPIPS curve overlaid on the figure.

\begin{figure}[h]
  \centering
  \includegraphics[width=\linewidth]{fig_convergence}
  \caption{Per-aspect mean human rating ($n=12$ participants, error
  bars are SEM) and offline LPIPS (teal dashed, axis inverted)
  across the four 3DGS checkpoints. Top row: \texttt{train}. Bottom
  row: \texttt{truck}. Columns are the three slider aspects.}
  \label{fig:convergence}
\end{figure}

These are descriptions of the data, not claims about NVS methods. A
researcher analyzing this data is the appropriate party to make
substantive claims; our role is to verify that the data has the
structure needed to support them. Appendix~\ref{app:analyses}
illustrates the kinds of analyses the data admits --- Spearman
correlations between A/B preference and offline-metric gap,
Mann--Whitney $U$ tests for aspect-citation effects, and (once
$n=18$ is large enough) Kendall's $W$ or ICC$(2,k)$ for inter-rater
agreement.

\subsection{Participant feedback}
\label{sec:participant_feedback}

This subsection addresses the participant-side success criterion via
two channels: indirect indicators from session logs, and a brief
post-session self-report survey administered to each participant.

\paragraph{Indirect indicators (from session logs).}
Every participant who loaded the viewer completed the session through
to submission: there were no mid-session abandonments in either case
study. Median session duration was approximately 9 minutes in case
study A and 5 minutes in case study B --- consistent with our piloting
estimate. No experimenter intervention was required
during any session, and no participant reported a stall, freeze, or
unresponsive UI in the post-submission notes field.

\paragraph{Post-session self-report ($n=9$).}
After each session we administer a five-question Likert survey
(1--5 scale) covering overall experience, mental fatigue, perceived
session length, instruction clarity, and ease of use of the
interactive viewer; the questionnaire itself is documented in
Appendix~\ref{app:satisfaction}. Table~\ref{tab:satisfaction} and
Figure~\ref{fig:satisfaction} summarize the $n=9$ survey responses
collected so far. The satisfaction survey was introduced partway
through recruitment so its sample lags behind the rating data;
numbers will be updated as additional survey responses come in.

\begin{table}[h]
\centering
\caption{Post-session participant survey, $n=9$ responses to date.
``Better'' indicates the direction toward positive participant
experience. Anchor conventions: overall (1 = poor, 5 = excellent);
fatigue (1 = none, 5 = a lot); length (1 = too short, 3 = just right,
5 = too long); clarity (1 = unclear, 5 = very clear); ease
(1 = hard, 5 = easy).}
\label{tab:satisfaction}
\small
\begin{tabular}{lcccl}
\toprule
Question & Better & Mean & Range & Distribution \\
\midrule
Overall experience  & $\uparrow$    & 4.44 & 4--5 & 4$\times$5, 5$\times$4 \\
Mental fatigue      & $\downarrow$  & 1.78 & 1--4 & 1$\times$4, 2$\times$4, 4$\times$1 \\
Session length feel & $\rightarrow$ & 3.11 & 3--4 & 3$\times$8, 4$\times$1 \\
Instruction clarity & $\uparrow$    & 4.44 & 4--5 & 4$\times$5, 5$\times$4 \\
Viewer ease of use  & $\uparrow$    & 3.56 & 2--5 & 2$\times$1, 3$\times$3, 4$\times$4, 5$\times$1 \\
\bottomrule
\end{tabular}
\end{table}

\begin{figure}[h]
  \centering
  \includegraphics[width=0.95\linewidth]{fig_participant_feedback}
  \caption{Distribution of post-session participant survey responses
  ($n=9$), one histogram per question. Bar heights are participant
  counts; the dotted red line on each panel marks the mean. The
  figure will be re-rendered as the survey sample grows.}
  \label{fig:satisfaction}
\end{figure}

The signal remains positive across the board. Overall experience
sits uniformly in the 4--5 band; instruction clarity is uniformly
$\geq 4$; and 8 of 9 participants rated the session length as
``just right'' (3), with one participant rating it 4
(``a bit too long''). Mental fatigue is low (1--2) for eight of nine
respondents, with one outlier at 4. The weakest dimension is ease of
use of the interactive viewer (mean 3.56). The 2 came from the same
participant whose row also shows elevated fatigue (4), suggesting
that for at least one participant the viewer UI was itself a source
of cognitive load rather than just the rating task; a new 5 in this
column from a recent respondent pulls the mean up. With $n=9$ we
report these patterns descriptively rather than test them, and will
revisit as additional survey responses arrive.

\subsection{Limitations}
\label{sec:limitations}

The most important limitation remains sample size. At $n=18$ and
$n=12$ the statistical tests above are reasonable for the per-scene
and per-stimulus signals we surface, but underpowered for finer-grained
between-aspect comparisons; we report effect sizes alongside $p$-values
where possible. Beyond sample
size, several threats to validity warrant explicit acknowledgment.

\emph{Display variance.} Participants connected over the public
internet from uncalibrated monitors. Differences in display brightness,
gamut, and refresh rate introduce noise we cannot correct for. The
case study A design --- both methods rendered to the same participant
in the same session --- partially controls for this, but case study B
(stimuli presented sequentially) does not.

\emph{Network variance.} WebRTC adapts bitrate dynamically based on
network conditions, which means a participant on a slower network may
be evaluating a different visual stimulus than one on a faster
connection. We do not currently log per-session bitrate; for a future
version this should be instrumented and either gated or stratified.

\emph{Selection bias.} Participants were recruited from the authors'
academic environment and self-selected for interest in 3D rendering;
generalization to the broader population is not claimed.

\emph{Order and learning effects.} Case study B fixes iteration order
as ascending within scene, trading order-effect contamination for a
stable reference frame for ``better''. Case study A randomizes scene
order and A/B side assignment per pair. A first-pair learning effect
may still inflate ratings on early scenes; the \texttt{bonsai}
practice trial is intended to absorb most of this but is not a
complete control.

\emph{Method-blinding integrity.} The viewer hides method identities
behind anonymous badges. We did not collect post-hoc data on whether
participants could guess which side was which, which would be the
standard check on blinding in a perception study.

\emph{Participant-side measurement is indirect.} As described in
\S\ref{sec:participant_feedback}, our current participant-side
success indicators are entirely indirect. The planned post-session
survey is the appropriate fix, but it is not yet in the field for the
current data.
